% !TEX program = xelatex
\documentclass[11pt,a4paper,fontset=none]{ctexart}
\input{preamble}
\expandafter\def\csname refnum@FUND-007\endcsname{1}
\expandafter\def\csname refnum@FUND-200\endcsname{2}
\expandafter\def\csname refnum@FUND-308\endcsname{3}
\expandafter\def\csname refnum@FUND-323\endcsname{4}
\expandafter\def\csname refnum@EXT-QA\endcsname{5}
\expandafter\def\csname refnum@FUND-400\endcsname{6}
\expandafter\def\csname refnum@FUND-300\endcsname{7}
\expandafter\def\csname refnum@FUND-220\endcsname{8}
\expandafter\def\csname refnum@FUND-215\endcsname{9}
\expandafter\def\csname refnum@FUND-356\endcsname{10}
\expandafter\def\csname refnum@FUND-251\endcsname{11}
\expandafter\def\csname refnum@FUND-359\endcsname{12}
\expandafter\def\csname refnum@FUND-001\endcsname{13}
\expandafter\def\csname refnum@FUND-316\endcsname{14}
\expandafter\def\csname refnum@FUND-113\endcsname{15}
\expandafter\def\csname refnum@FUND-114\endcsname{16}
\expandafter\def\csname refnum@FUND-115\endcsname{17}
\expandafter\def\csname refnum@FUND-116\endcsname{18}
\expandafter\def\csname refnum@REG-010\endcsname{19}
\expandafter\def\csname refnum@FUND-119\endcsname{20}
\expandafter\def\csname refnum@EXT-QC\endcsname{21}
\expandafter\def\csname refnum@EXT-AB\endcsname{22}
\expandafter\def\csname refnum@REG-001\endcsname{23}
\expandafter\def\csname refnum@REG-005\endcsname{24}
\expandafter\def\csname refnum@REG-002\endcsname{25}
\expandafter\def\csname refnum@REG-003\endcsname{26}
\expandafter\def\csname refnum@REG-004\endcsname{27}
\expandafter\def\csname refnum@REG-014\endcsname{28}
\expandafter\def\csname refnum@REG-006\endcsname{29}
\expandafter\def\csname refnum@REG-009\endcsname{30}
\expandafter\def\csname refnum@REG-012\endcsname{31}
\expandafter\def\csname refnum@REG-007\endcsname{32}
\expandafter\def\csname refnum@REG-008\endcsname{33}
\expandafter\def\csname refnum@FUND-004\endcsname{34}
\expandafter\def\csname refnum@FUND-009\endcsname{35}
\expandafter\def\csname refnum@FUND-304\endcsname{36}
\expandafter\def\csname refnum@FUND-403\endcsname{37}
\expandafter\def\csname refnum@FUND-352\endcsname{38}
\expandafter\def\csname refnum@FUND-312\endcsname{39}
\expandafter\def\csname refnum@FUND-150\endcsname{40}
\expandafter\def\csname refnum@FUND-011\endcsname{41}
\expandafter\def\csname refnum@FUND-160\endcsname{42}
\expandafter\def\csname refnum@FUND-013\endcsname{43}
\expandafter\def\csname refnum@REG-013\endcsname{44}
\expandafter\def\csname refnum@REG-015\endcsname{45}
\expandafter\def\csname refnum@REG-016\endcsname{46}

\begin{document}
\hypersetup{pageanchor=false}
\begin{titlepage}
\thispagestyle{empty}\setlength{\parindent}{0pt}
\vspace*{18mm}
{\sffamily\large\color{teal}MSIM CHINA · INDUSTRY RESEARCH\par}
\vspace{18mm}
{\Huge\bfseries\color{navy}基金公司LLM应用调研\par}
\vspace{8mm}
{\Large\color{muted}模型任务、输入输出与实施条件\par}
\vspace{15mm}\textcolor{teal}{\rule{\linewidth}{2pt}}\vspace{9mm}
\begin{tabularx}{\linewidth}{@{}P{28mm}Y@{}}
研究主题 & 投研分析、数据查询、知识问答与合规审核、LLM辅助量化\\
研究范围 & 中国公募基金，以及海外资管与公开研究中的相关案例\\
资料截止 & 2026年9月9日\\
项目阶段 & 本地LLM部署项目：Phase 1行业研究\\
\end{tabularx}
\vfill
{\large\color{navy}明确LLM在基金业务中的具体任务，\\[4pt]为本地部署与应用验证提供依据。\par}
\vspace{15mm}
{\small\color{muted}2026年9月\par}
\end{titlepage}
\hypersetup{pageanchor=true}

\section*{摘要}
\addcontentsline{toc}{section}{摘要}
基金公司已在投研分析、数据查询、知识问答和合规审核中使用LLM处理文本与业务指令，并在量化研究中探索事件提取、因子构造和研究工具调用。公开案例显示，模型承担的工作可以落实为具体产物：研报小结、研究问题分解、查询语句、条款答复、风险提示，以及因子表达式和修改建议。模型产物通常还要交由数据库、计算程序或业务人员继续处理。\src{FUND-007}\src{FUND-300}\src{FUND-200}\src{FUND-308}

投研与知识应用中，LLM读取研究问题及检索到的材料，提取信息并组织答案。南方交易助理进一步将模型输出用于执行：LLM识别意图和参数，生成工作流或SQL，调度引擎与数据库执行后，再由回复环节生成面向使用者的解释。模型因此既可以生产分析文本，也可以生产供程序执行的指令；权限检查与数值计算则由相应系统承担。\src{FUND-200}\src{FUND-220}

量化应用需要区分生成研究方案与计算研究结果。国内材料已披露LLM从舆情中提取政策敏感度因子、辅助评价既有因子等任务，但部分输出格式仍未公开。QuantaAlpha的公开实现提供了更具体的模型调用证据：LLM将研究方向展开为候选假设，生成或修正因子表达式，并根据程序返回的检验结果提出下一轮方向。回测收益、风险指标和因子数值由程序计算。Man AHL的AlphaTrend则展示了模型生成信号方案、参与实现并综合回测证据形成研究结论的分工。\src{FUND-007}\src{EXT-QA}\src{EXT-QC}\src{FUND-323}

同行实践已形成可借鉴的方法，开展受控验证有助于积累模型与业务系统配合的经验。本地部署评估应首先明确模型需要接收的材料、输出的格式及后续使用方式，再比较模型质量和资源要求。数据使用、外部处理、系统权限及模型管理同时受到法律和内部制度约束，具体范围需由相应责任方确认。现阶段适合从资料明确、输出可核验的内部任务切入，逐步验证其适用性。\src{REG-002}\src{REG-001}\src{REG-003}
\clearpage
\tableofcontents
\vfill
\smallnote{公司材料以2026年公开披露为主，2025年MENTOR列为方法补充。输入、任务与产物按来源中的实际说明整理；未公开的接口、提示词和字段格式不作推定。文中“输出示意”为解释任务而构造，不是公司真实运行记录。公司效果和论文结果均未在本项目独立复现。}
\clearpage

\section{行业概览与模型职责}
本次调研以2026年二季度末非货公募规模前20家公司为主要样本，并保留大成的补充案例。公司规模、模型任务、主要产物和技术合作信息见\hyperref[tab:landscape]{表A（附录A）}。规模剔除ETF联接基金，来自财联社报道，主要样本约为3,448亿元至1.70万亿元。该口径用于界定调研对象，不用于推断技术能力或估计全行业采用率。\src{FUND-400}

已披露的应用分布在多个业务环节。富国使用本地LLM生成研报小结和市场观点，天弘的FinAgent结合模型与工具组织财务分析，南方使用LLM解析查询需求并生成SQL。鹏华通过模型提取资讯标签、生成知识答复和审核建议，汇添富将模型用于托管风险邮件提取与风控条款问答。永赢的邮件反钓鱼应用则说明LLM也进入了运营安全场景。\src{FUND-007}\src{FUND-300}\src{FUND-200}\src{FUND-308}\src{FUND-220}\src{FUND-352}

应用名称不足以说明模型的作用。例如，“智能投研”可能指文本归纳，也可能指生成研究计划；“因子挖掘”可能由LLM提出表达式，也可能主要由传统机器学习完成。表\ref{tab:roles}将四个方向按模型接收的材料和直接产物区分。判断LLM贡献时，需要定位这些产物，再检查其是否被后续流程采用。

\begin{table}[H]\small\centering
\caption{四个方向中的LLM任务与直接产物}\label{tab:roles}
\begin{tabularx}{\linewidth}{@{}P{23mm}Y Y Y@{}}\toprule
方向 & 输入模型的材料 & LLM承担的工作 & 模型直接产物\\\midrule
投研分析 & 研究问题、相关研报或财报、工具返回数据。 & 归纳证据，拆解问题，组织解释或预测。 & 分项小结、分析框架、预测清单或排序。\\
数据查询 & 自然语言需求、表字段说明或接口目录。 & 识别意图与参数，选择数据操作，生成查询。 & 查询计划、接口参数或SQL；执行后生成答复。\\
知识与合规 & 问题或待审材料、检索到的依据及规则。 & 提取信息，解释条款，识别疑点并建议修改。 & 带依据的答案、资讯标签、风险信息或修改建议。\\
LLM辅助量化 & 文本材料、研究方向、变量与算子说明、检验反馈。 & 提取事件，提出假设，构造表达式，解释反馈并修订。 & 事件记录、研究方向、因子表达式、修正建议。\\\bottomrule
\end{tabularx}
\source{按正文案例归纳。表中列示任务层面的材料和产物，不代表各公司采用统一接口或输出格式。}
\end{table}

表中产物的使用方式也不同。研报小结主要供研究员阅读，SQL会进入数据库执行，因子表达式会进入计算和回测。涉及执行时，模型提出的内容还需通过权限、格式或程序检查。记录检索、调度和回测各自的处理结果，有助于识别模型对整个流程的实际贡献。

公开资料对这些分工的说明程度并不一致。南方披露了较具体的查询链条，QuantaAlpha公开了部分提示词和代码；部分基金仅介绍业务功能，尚无模型调用记录或输出样本。技术名称也不等于完整合作关系：天弘列出的Wind和同花顺说明数据工具来源，工银瑞信的历史FundGPT材料则明确列出恒生聚源、工银科技和智谱AI等合作方。两者分别提供工具与合作证据，实施和运维分工仍以实际披露为限。\src{FUND-200}\src{EXT-QC}\src{FUND-300}\src{FUND-316}
\clearpage

\section{投研分析：LLM如何组织研究材料}
\subsection{富国：将研报内容归纳为分项小结与市场观点}
富国的市场观点流程使用本地LLM、提示词和研究数据服务。研究员指定投资标的与分析方向后，系统先依据研报与标的的相关程度筛选材料。LLM在该流程中承担内容分析和生成：对选取的研报形成知识小结，按指定方向归纳，再综合小结生成市场观点。生成的小结与综合观点供研究员继续分析和复核。\src{FUND-007}（PDF第8--9页）

例如，研究方向若为某行业的需求变化，输入模型的材料可以是筛选后的研报内容及分析要求，输出应是各材料的需求判断和综合观点。公开材料尚未说明相关性评分的具体实现，也未提供完整提示词和内部字段。

富国另有知识问答流程：资料系统对研报、行情、内部点评等内容进行切片、索引和存储，检索组件找出相关片段后，LLM根据这些片段总结并回答问题。模型输出自然语言答案，系统在回答中标记数据来源。检索增强生成（RAG）由此将资料检索与答案生成连接起来：检索组件提供材料，LLM据此组织回答。\src{FUND-007}（PDF第9页）

公司报告，研究系统包含约30万份研报，相关流程上线后，原来一周以上的工作可缩短至3小时以内。该数字反映整套流程的自报效果；原文没有给出统一任务样本和人工复核工时。后续验证可先检查模型小结是否保留关键事实、综合观点是否得到原文支持，再衡量整理与复核的总耗时。

\subsection{天弘：LLM参与问题拆解与结果组织}
天弘FinAgent的公开架构将模型、工具与业务智能体连接。系统接收复杂研究任务后，将其拆为子任务，调用金融数据接口或计算工具，再综合各环节的结果。LLM在其中承担任务理解、分析内容组织和语言生成，工具负责数据获取、指标计算等操作。公司尚未公开完整节点提示词和调用记录，因此以下按已披露业务流程说明模型职责，不推定每个节点的内部实现。\src{FUND-300}（PDF第4--6页）

盈利归因案例能够具体说明这种分工。输入是研究员关于企业盈利驱动的指令，以及财报和后续取得的数据；FinAgent将净利润问题展开为收入、成本、费用，并进一步整理销量、单价、原材料和渠道等因素。模型参与形成的主要产物是结构化归因分析框架和解释文本，供研究员判断增长质量与持续性。财务数值及其计算结果需要来自财报、数据接口和金融计算库。\src{FUND-300}（PDF第6--7页）

一份可核验的输出因此至少应能区分三类内容：材料已经说明的事实、程序已经计算的变化，以及仍待验证的解释。公司报告相关手工整理和计算由数日缩短至分钟级，但未披露归因准确率或统一对照。衡量模型是否减少研究准备工作，需要同时检查分析框架的质量与复核成本。

\subsection{MENTOR：多次LLM调用产生预测和策略文本}
MENTOR是易方达与高校共同作者的2025年研究。与单次研报摘要不同，它将LLM用于多个相连任务，每次调用产生供下一步使用的文本或列表。官方补充材料公开了事件预测、行业排序和反馈提示词，能够比一般系统介绍更清楚地定位模型输入与输出。\src{FUND-113}\src{FUND-114}

事件预测调用读取上周热点和改进指导，生成下一周按市场影响排序的10条事件；行业排序调用读取事件、历史表现和策略，返回行业名称的有序列表。周期结束后，反馈调用比较预测与实际结果，生成用于后续预测的改进意见。该机制更新提示词中的策略文本，模型权重保持不变。\src{FUND-114}（PDF第2--3页）

\begin{table}[H]\small\centering
\caption{MENTOR中不同LLM调用的产物}
\begin{tabularx}{\linewidth}{@{}P{28mm}Y Y@{}}\toprule
模型任务 & 主要输入 & 直接输出及用途\\\midrule
预测事件 & 上周事件、改进指导。 & 下一周事件描述清单，进入后续排序与评价。\\
排列行业 & 事件信息、历史行业表现、排序指导。 & 行业名称的有序列表，供程序计算排序指标。\\
诊断偏差 & 先前预测、实际事件或排序结果。 & 改进意见和策略文字，用于后续周期的提示。\\\bottomrule
\end{tabularx}
\source{MENTOR官方补充材料第3节。列表和指导文本是模型产物；排序相关性等统计指标由评价流程计算。}
\end{table}

原作者实验显示，MENTOR在部分行业排序指标上优于所列基线，但模型版本之间存在取舍，具体结果见\hyperref[tab:mentor]{附录C表\ref*{tab:mentor}}。对模型应用的直接启示是，中间产物可以逐项检验：事件清单是否准确、行业列表是否完整、修改后的策略是否在后续数据上有效。只检查最终回测收益，无法识别错误来自信息提取、预测还是排序。该案例属于学术实验，完整数据的可得性和时间污染仍需在复现前核查。

\Needspace{7\baselineskip}
\section{数据查询：LLM如何把需求转为可执行指令}
\subsection{南方：意图、参数、SQL与结果答复}
南方“小喃同学”的公开方案明确将微调大模型用于意图识别和参数解析。模型接收自然语言指令，将其拆为子任务并编排工作流；调度引擎执行该工作流，在节点上调用API、SQL查询、知识检索或Python解释器。模型首先产生的是对需求的结构化理解及任务安排，执行结果由各工具返回。\src{FUND-200}（PDF第4页）

在SQL子流程中，表格选择智能体接收问题和数据库元数据，包括表名称、业务说明、字段注释及常见查询方式，形成候选表。SQL撰写智能体再读取选中的表结构和查询模板，生成包含筛选、关联或分组逻辑的查询语句。LLM的核心工作是将业务语言转换为数据操作表达，数据库仍负责执行语句，查询引擎负责语法、执行计划和分页等检查。\src{FUND-200}（PDF第5--6页）

\begin{table}[H]\small\centering
\caption{南方查询流程中的模型输出与执行分工}
\begin{tabularx}{\linewidth}{@{}P{27mm}Y Y Y@{}}\toprule
环节 & 模型接收的信息 & 模型或智能体产物 & 后续处理\\\midrule
解析需求 & 自然语言问题及上下文。 & 意图、参数、子任务与工作流。 & 调度引擎选择执行节点。\\
选择数据 & 问题、表与字段说明。 & 候选表及与问题对应的数据范围。 & 向查询生成环节提供表结构。\\
生成查询 & 表结构、问题、SQL示例模板。 & SQL查询语句。 & 程序检查权限与语法，数据库执行。\\
解释结果 & 查询结果和原始问题。 & 与问题对应的自然语言答复。 & 统计和图表由相应计算功能支持，使用者复核口径。\\\bottomrule
\end{tabularx}
\source{南方案例PDF第4--6页。来源未公开各节点的完整消息格式，不据此还原内部接口。}
\end{table}

若任务是按行业汇总指定评级和期限债券的余额，模型需要从问题中识别统计日、评级、期限与分组要求，再生成相应查询。这是对模型任务的说明性例子。模型的直接产物可以是SQL和结果解释；余额求和来自数据库执行，访问范围由系统控制。当前评级替代历史评级、重复关联导致重复统计等错误，都需要通过实际查询和返回结果核验。

南方报告整套交易助理整合5个以上系统及20个以上流程后，日均操作耗时下降约40\%。该指标未单独衡量LLM生成SQL的准确性。验证此类应用时，需要分别记录意图与参数是否正确、查询是否符合业务口径、执行是否受控，以及答复是否忠实于返回数据。

\subsection{易方达Index-Hub：模型选择接口并组织数据答复}
Index-Hub的公开客户端采用预定义接口。其任务说明要求使用者所用的AI智能体先识别查询对象，再依据接口目录选择查询能力，通过客户端取得真实数据，最后组织比较和解释。LLM主要输出接口选择、查询参数和自然语言答复；客户端承担鉴权、缓存和请求，数据服务返回实际数值。\src{FUND-115}\src{FUND-116}

以“查询沪深300前十大成分股”为输出示意，模型需要先识别对象为指数、代码为000300，再选择成分股查询能力。待接口返回成分股及权重后，模型形成带数据日期的说明。回答中的数值以接口返回为依据。该例根据公开任务说明构造，尚未实际执行。

与南方的SQL生成相比，Index-Hub将可执行操作预先限定在接口目录中，模型需要作出的选择更集中。公开代码足以研究这部分任务组织，却没有公开易方达内部模型或数据后端；固定版本的部分子包和授权条件仍需核实。因此，可迁移的是模型与工具之间的分工，实际部署还需验证接口可用性和任务质量。
\Needspace{7\baselineskip}

\section{知识与合规：LLM如何生成答案和审核意见}
\subsection{鹏华：检索结果是输入，答案与标签是产物}
鹏华的知识库方案明确区分数据工程和模型生成。知识系统先解析、清洗和切分资料，通过嵌入模型转换为向量后入库。收到问题后，检索组件取得相关片段，再将片段与原问题提交给LLM。生成式LLM执行上下文理解和信息整合，直接输出基于资料的自然语言答案。检索向量的生成和相似度搜索属于另外的组件，不宜与答案生成混为一次模型推理。\src{FUND-308}（PDF第4页）

投研资讯助手还将LLM用于信息抽取和聚合。输入包括公告、新闻和券商研报；模型提炼关键信息并生成可定制标签，再根据这些结构化信息组织个股舆情要点和摘要。其产物至少包括信息标签和摘要报告，供研究员筛选事件与阅读。完整标签定义、逐条输出格式和独立准确率尚未公开。\src{FUND-308}（PDF第5--6页）

\subsection{鹏华与汇添富：从理解材料到生成待复核事项}
鹏华营销审核使用大模型模拟合规审核思路，并结合7类、71项规则检查风险提示、基金信息、业绩宣传和禁止性表达。模型相关产物是风险识别结果和修改建议，供合规人员复核。规则是否适用、业绩数值是否真实、材料是否可以发布，仍涉及规则维护、事实核对和业务审批。各项规则在模型与程序之间的具体分工，以及完整提示词和审核接口，尚未公开。\src{FUND-308}（PDF第6页）

公司报告累计审核1,266篇文件、执行近6万项信息审核，产生989条有效修改建议，最新采纳率91\%。这些数量支持系统已有使用和建议产出的判断，但采纳率不等于审核准确率，尚需违规标注和漏检数据才能评价模型判断质量。

汇添富将大模型用于托管行风控邮件自动化提取，以及风控条款智能问答。前一任务把邮件文字转为可供工作台处理的信息，后一任务结合风控规则与法规、内控或合同原文的绑定关系，帮助定位拦截背后的依据。模型输出的是提取信息和条款解释，既有风控规则引擎继续执行限额等计算。原文没有公开邮件字段清单，也未将合同规则自动填充列为已完成的准确性验证。\src{FUND-220}（PDF第5、10--11页）

\begin{table}[H]\small\centering
\caption{知识与合规应用中的输入、模型动作和产物}
\begin{tabularx}{\linewidth}{@{}P{25mm}Y Y Y@{}}\toprule
任务 & 模型输入 & LLM承担的动作 & 产物与后续使用\\\midrule
知识问答 & 问题、检索片段、来源信息。 & 理解依据并组织回答。 & 答案及来源说明，供使用者核对。\\
资讯提取 & 公告、新闻或研报内容及标签要求。 & 识别关键信息、生成标签、归纳内容。 & 资讯标签、个股要点及摘要，进入研究阅读与筛选。\\
营销审核 & 待审内容与相关审核要求；完整内部上下文未公开。 & 理解表达，识别疑点并建议修改。 & 待复核建议，进入合规审核。\\
风险邮件与条款 & 邮件文字，或问题及对应规则原文。 & 提取风险信息、解释条款依据。 & 风控工作台信息或条款答复，后续操作由业务系统和人员处理。\\\bottomrule
\end{tabularx}
\source{鹏华和汇添富公开案例。按披露整理产物类别，不补写未公开字段。}
\end{table}

例如，待审文案出现“适合所有投资者”时，模型可以在相关规则上下文中识别适用范围疑点，并输出问题片段和修改建议。这是输出形式的说明性例子，不代表来源公布过该条真实审核记录。复核人员需要检查规则是否适用，以及修改建议是否保留了产品事实。

这些应用的共性是，模型先将文字转为答案、标签或待处理事项，再由人员或系统继续处理。量化研究也采用类似分工，但其产物还包括可进入计算流程的事件记录、表达式和程序，因而需要进一步说明每次模型调用的职责。
\Needspace{7\baselineskip}
\section{LLM辅助量化：模型生成什么，程序计算什么}
\subsection{文本信号：模型将自然语言转为事件信息}
国内基金已披露的量化应用，首先体现为从非结构化文本中提取研究信息。富国明确使用LLM从舆情提取政策敏感度因子，并将既有因子的绩效分析数据输入LLM，获得辅助评价和参考打分。前者的模型输入是舆情文本，产物被称为政策敏感度因子；后者输入程序已经计算的绩效数据，产物是评价意见和参考分。原文尚未说明政策因子的具体取值与计算关系，也未披露评价分数的范围和逐条输出格式。\src{FUND-007}（PDF第8页）

广发的媒体调研提到利用大模型从非结构化数据挖掘另类因子，华泰柏瑞介绍了公告、纪要和舆情结构化后与量化信号结合的应用。两者能够说明LLM参与文本处理，但仍缺少具体提示词、字段定义与模型输出样本。华泰柏瑞有关自主执行数据、代码和回测的Agent属于展望，不列作已经公开验证的完整研究流程。\src{FUND-251}\src{FUND-359}

以下扩产公告仅用于示意模型产物。假设公告写明“拟将年产能由10万吨提高至15万吨，预计两年后投产，尚待审批”，任务要求LLM读取原文并提取事实，模型可以输出：

\begin{quote}\small
\textbf{事件记录示意}\par
事件：扩产；状态：计划、待审批；原产能：10；计划产能：15；单位：万吨；预期投产：两年后；证据：对应公告原句。
\end{quote}

这里LLM承担的是实体与事件识别、数量和单位提取，以及对“拟”“尚待审批”等语义的区分。程序随后计算计划增幅、匹配证券、去重并附加首次可得时间，再将记录接入研究数据集。例如，50\%的计划增幅由程序根据提取值计算，研究中的收益和风险也由后续统计检验获得。原文证据与发布时间同时保留，才能检查模型是否遗漏了状态限制或混入后续信息。

因此，该方向至少有两项独立的验收对象：LLM生成的事件记录是否准确，以及该记录进入研究后是否有增量价值。前一项可以用人工标注评价字段和证据，后一项需要样本外研究。将两者分开，才能判断问题来自语言理解还是研究假设本身。

\subsection{QuantaAlpha：LLM生成方向、表达式和修订意见}
QuantaAlpha为2026年的外部研究框架，不属于已确认的国内基金生产系统。其方法的关键是多次调用LLM，每次调用处理不同上下文，并将结果传给下一环节。公开论文与固定代码版本能够定位到研究方向规划、因子构造、一致性核对和方向修改等任务。\src{EXT-QA}\src{EXT-QC}

\paragraph{研究方向规划。}
规划提示词接收初始研究方向和需要展开的方向数量，要求LLM生成具体、可检验且具有差异的探索方向。模型返回包含\texttt{directions}数组的JSON，框架解析数组后分派研究任务。这里模型生成的是自然语言研究命题，尚未生成因子数值。代码在禁用模型或解析失败时可以采用预设方向，因此实际复现还需区分模型返回与程序回退。\src{EXT-QC}（\texttt{planning.py}及规划提示词）

\paragraph{假设与因子构造。}
论文中的想法智能体根据市场观察、金融知识、参数要求或已有研究记录提出假设；因子构造环节进一步将假设转成语义说明、数学定义和由允许算子组成的符号表达式。LLM在这里决定候选机制如何表述、采用哪些字段和操作，以及如何组合为一个因子定义。这些规格随后交给代码组件实现，并在历史数据上计算因子序列。\src{EXT-QA}（第4--5页）

论文规定，验证后的表达式经代码组件转换为可执行实现；当实现失败时，可以触发基于LLM的修复。模型读取失败信息及目标定义，提出修改后的表达式或实现，程序再执行检验。因子值、组合构造和回测统计由程序运行产生。代码中的\texttt{factor\_calculate}与\texttt{factor\_backtest}分别调用开发和运行组件，体现了这一分工。\src{EXT-QA}\src{EXT-QC}（\texttt{loop.py}）

\Needspace{12\baselineskip}
\paragraph{一致性核对。}
公开提示词将假设、因子说明、数学公式、符号表达式和变量说明一并提交给LLM，要求模型判断各层含义是否一致。模型返回的JSON包含一致性判断、问题严重程度、反馈及可选的修正表达式。其核心输出形式可概括为：

\begin{quote}\small\ttfamily
is\_consistent: 一致性判断\par
severity: 问题严重程度\par
 overall\_feedback: 不一致之处及修改意见\par
 corrected\_expression: 建议的修正表达式（如有）
\end{quote}

上述字段取自公开提示词，中文说明概括其含义。模型提供的是对语义关系的判断，框架解析后才决定是否重试或采用修正。原提示词允许一定窗口和归一化差异，代码还存在异常时跳过检查的分支，实际使用时还需独立进行数值、语法和时序验证。\src{EXT-QC}（\texttt{consistency\_prompts.yaml}、\texttt{consistency\_checker.py}）

\paragraph{根据反馈修改研究方向。}
修改调用接收父研究记录中的假设、表达式、回测指标和评价意见，LLM输出新假设、探索方向、差异理由及预期特征。交叉调用则读取多条记录，生成融合假设和组合理由。指标由回测器提供，LLM根据这些指标撰写诊断和后续研究建议。提示词要求方向“正交”或“互补”，只是生成要求；实际相关性仍需由程序计算。\src{EXT-QC}（演化提示词及\texttt{mutation.py}）

\begin{table}[H]\small\centering
\caption{QuantaAlpha的模型调用与后续处理}
\begin{tabularx}{\linewidth}{@{}P{24mm}Y Y Y@{}}\toprule
模型调用 & 输入上下文 & LLM直接输出 & 后续程序\\\midrule
规划方向 & 初始方向、数量及差异要求。 & 研究方向数组。 & 解析JSON、建立研究分支。\\
构造因子 & 假设、字段与算子约束、已有记录。 & 因子说明、公式及符号表达式。 & 解析表达式，构建实现并计算因子。\\
核对与修正 & 假设、说明、表达式及问题反馈。 & 一致性意见、修正表达式或实现。 & 根据配置重试，并执行程序验证。\\
解释结果 & 研究命题、程序产生的指标与记录。 & 原因分析、保留或修改意见。 & 保存反馈并提供给后续调用。\\
修改或组合 & 一条或多条父记录及其反馈。 & 新假设、探索方向或融合理由。 & 开始下一轮构造与检验。\\\bottomrule
\end{tabularx}
\source{论文第4--7页与固定代码版本。表述限于已取得的方法和选取文件，不表示完整运行已复现。}
\end{table}

这一过程可以用图\ref{fig:llm-program}概括。图中的往返信息明确了模型与计算环境的联系：LLM产生可执行规格和修订文本，程序产生执行结果与指标，两类输出通过框架连接。

\begin{figure}[H]\centering
\begin{tikzpicture}[node distance=10mm]
\node[flowwide,fill=teal!8,text width=39mm,minimum height=20mm](a){LLM生成\par 假设、因子说明\par 与表达式};
\node[flowwide,right=of a,fill=gray!8,text width=39mm,minimum height=20mm](b){程序执行\par 实现、计算\par 与回测};
\node[flowwide,right=of b,fill=teal!8,text width=39mm,minimum height=20mm](c){LLM读取反馈\par 生成诊断与\par 修订建议};
\draw[arr](a)--node[above,font=\scriptsize]{规格}(b);
\draw[arr](b)--node[above,font=\scriptsize]{结果}(c);
\draw[darr](c.south)--++(0,-.65)-|node[pos=.25,below,font=\small]{新的研究方向与修改要求}(a.south);
\end{tikzpicture}
\caption{依据QuantaAlpha公开材料归纳的模型产物与程序结果传递；非原始运行日志。}\label{fig:llm-program}
\end{figure}

\subsection{一个表达式如何进入下一轮模型调用}
QuantaAlpha附录C给出了一条具体研究记录。系统希望组合持续动量与脆弱动量，并按市场波动状态调整权重。生成的子因子实际使用20日价格与成交量变动的相关性，乘以5日日内收益均值，再作横截面排名。模型相关产物在这里是因子假设与表达式，随后才由程序计算对应的历史序列和回测结果。\src{EXT-QA}（附录C，第17--20页）

程序返回的结果显示，子因子的Rank IC高于父记录，但信息比率略降、最大回撤扩大。原作者保存的反馈据此拒绝直接采用，并建议加入波动状态后重新检验。对LLM作用的观察应落在这一步：模型读取相互不一致的结果，输出解释和修改方向，框架再将这些文字传入下一轮。原文未报告该次修改完成后的验证结果。

该记录同时暴露了需要人工复核的语义偏差。假设提及的机构行为、投机活动和波动权重，并未全部体现在价量表达式中。LLM给出的经济解释可能比实现内容更丰富，需要检查表达式是否真正测量了所述对象。详细公式口径、案例数字及总体实验结果列于附录C。

原作者总体实验报告了同一基础模型下的框架差异，并通过消融实验考察修改环节的贡献，见\hyperref[tab:qa]{附录C表\ref*{tab:qa}}。这些结果为进一步研究LLM生成与迭代提供依据；具体解释是否成立，以及本地模型能否达到相近质量，仍需分别验证。

\subsection{Man AlphaTrend：模型生成方案、参与实现并撰写研究结论}
Man AHL的AlphaTrend将研究流程预先划分为想法生成、实现、回测、回测分析和综合结论等节点。LLM可以在多个节点被调用，后续调用使用先前步骤的结果；多次并行请求用于产生不同的信号方案，并比较输出是否一致。框架预先设定工作流顺序和工具依赖，LLM在各节点处理具体研究任务。\src{FUND-323}

在想法生成阶段，输入是研究员给出的趋势研究命题和任务上下文，模型输出候选信号方案或参数变体。在实现阶段，候选方案进入代码实现环节；在分析阶段，模型利用回测和基准分析结果形成解释。最终产物是关于方案表现、差异和后续研究价值的综合结论。原文展示了节点分工，但未公开各节点的完整提示词、代码和消息格式。

\begin{table}[H]\small\centering
\caption{AlphaTrend中模型产物的变化}
\begin{tabularx}{\linewidth}{@{}P{27mm}Y Y@{}}\toprule
阶段 & 提供给模型的内容 & 模型参与形成的产物\\\midrule
提出方案 & 原始研究命题、约束和上下文。 & 多个候选信号说明及实现思路。\\
落实方案 & 候选定义和实现任务。 & 供回测工具使用的候选实现；完整代码未公开。\\
分析与综合 & 执行结果、回测指标、基准及相关性比较。 & 对有效、较差或不确定方案的解释与研究结论。\\\bottomrule
\end{tabularx}
\source{Man原文工作流及“Three experiments”。市场数据和回测统计来自工具，表格不表示公开了完整模型消息记录。}
\end{table}

作者用已知可改进、预期较差和开放性想法测试该流程。恢复被移除机制后，方案表现改善；较差方案和开放方案则分别产生负面或有条件的结论，具体结果见附录C。该实验考察模型能否根据返回结果形成与证据相符的判断。结果来自历史模拟，完整实现和统一样本区间尚未公开。

\subsection{研究工具调用与代码生成}
博时BSBox和中欧的2026年材料介绍了数据访问、研究工具、Skills和运行环境，但尚未提供完整因子研究任务的输入输出记录。目前可确认的是平台提供了研究工具；具体因子任务中模型如何选择和使用这些工具，仍缺少可核验的完整记录。\src{FUND-215}\src{FUND-356}

易方达参与的NLPCC 2026投资Agent任务则提供了可读研究基线。LLM使用新闻、行情及任务约束形成投资决策，由回测服务执行模拟并返回结果。该案例属于学术任务，仍有隐藏评测和数据使用条件。\src{FUND-119}

代码生成任务可以更明确地观察LLM产物。AlphaQT-Bench将量化任务说明、输入数据结构和输出要求提交给模型，模型返回实现因子的Python代码。测试程序随后运行代码，与专家参考实现比较，并检查时间因果性和计算结构。LLM负责把自然语言要求转成程序，专家代码、测试数据和验收规则承担独立评价。\src{EXT-AB}（任务示例、第3.4节及末页提示词）

例如，时间截断检查先用包含较晚数据的全样本计算，再截断至较早时点重算，比较相同历史时点的输出。测试程序据此检查代码是否因较晚数据而改变历史输出。原论文中的验证正确率还同时要求可执行、功能正确及结构合规，具体结果见附录C。代码可运行只是模型产物的一项性质，不代表其已正确实现金融任务。

\subsection{本地模型替换应验证相同产物}
上述框架允许将模型服务与计算工具分开配置。QuantaAlpha通过\texttt{APIBackend}调用LLM，其代码中的\texttt{use\_local}控制本地或容器中的回测，并不表示LLM已在本地运行。原作者采用托管模型API和CPU回测，模型运行位置需要另行判断。\src{EXT-QC}（\texttt{loop.py}、模型客户端）

因此，评估本地部署时，应固定任务、输入材料和工具，检查本地模型能否生成同样可用的产物：研究方向能否解析、SQL或表达式是否符合约束、修正是否解决实际错误、研究结论是否忠实于程序结果。随后再衡量响应时间、资源和人工复核成本。接口能够连接，只证明服务可以调用；模型产物的质量仍需逐项验证。
\Needspace{7\baselineskip}

\section{合规要求与模型产物的使用}
LLM接收的材料和输出的用途决定了需要核查的管理条件。研报摘要涉及资料使用权与团队隔离，生成SQL涉及数据访问和执行权限，审核建议涉及规则版本与发布责任，研究代码则涉及运行环境及模型风险管理。法律与制度要求应结合这些具体处理活动逐项确认。

《生成式人工智能服务管理暂行办法》第2条以向境内公众提供生成式AI服务为适用条件，并明确排除未向公众提供此类服务的组织研发和应用。内部研究工具仍需遵守相应的数据、安全和金融机构管理要求。任务改为客户问答或对外发布时，服务对象和内容责任需要重新确认。\src{REG-001}\src{REG-005}

《个人信息保护法》要求确定处理依据，并对委托处理、敏感信息、自动化决策和出境等活动规定相应义务；《数据安全法》还涉及个人信息之外的数据分类与保护。对模型输入而言，需要确认哪些资料允许使用、由谁访问及如何保存；对输出而言，需要避免将受限原文、检索片段或查询结果提供给无权限使用者。相应权限需要由检索和执行系统实际控制。\src{REG-002}\src{REG-014}\src{REG-010}

\begin{table}[H]\small\centering
\caption{模型输入输出对应的控制事项}
\begin{tabularx}{\linewidth}{@{}P{27mm}Y Y@{}}\toprule
模型产物 & 主要使用风险 & 实施时需要的控制证据\\\midrule
摘要与知识答复 & 引用受限、过期或不足以支持结论的材料。 & 资料使用权、版本与检索权限，以及答案对应的原文。\\
SQL与工具参数 & 查询越权，或指令超出任务范围。 & 服务端最小权限、实际请求与返回值、执行前检查。\\
审核或风险意见 & 漏检、错误规则解释，或建议被直接当作批准。 & 有效规则、人工标签、复核意见及发布决定。\\
因子规格与代码 & 实现偏离假设，读取未授权数据，或执行不受控。 & 代码与数据版本、隔离环境、程序测试和采用记录。\\\bottomrule
\end{tabularx}
\source{表中为基于公开规则与案例提出的控制安排；具体义务和条件见附录B。}
\end{table}

外部OCR、嵌入服务、模型API、日志和远程运维都可能参与处理资料。本地生成模型并不自动消除这些外部环节，API也不必然形成出境。应核对实际处理位置、供应商角色、数据类型和相关数量，再判断适用要求。这些条件应在涉及外部处理的原型阶段确认。\src{REG-004}\src{REG-006}

《证券基金经营机构信息技术管理办法》第43条要求机构保留自身责任与重要系统控制，并对外部独立实施重要系统运维和日常安全管理作出限制，保留证监会另有规定的例外；第44--45条涉及审查、相应报送、协议和替换安排。数据服务、软件开发与运行维护的职责应分别明确。\src{REG-003}

模型风险管理还需覆盖生成内容的用途、验证和变更。AMAC《基金经营机构大模型技术应用规范》提供场景评测、权限及解释参考，属于团体标准，具体采用依据需要确认，资料性附录不能直接当作统一强制要求。美国SR 26-2及配套MRM指引面向所述银行机构，且其模型定义不涵盖生成式及Agentic AI，不能替代MSIM内部制度。内部政策、用例分级和COD允许的技术条件目前尚未取得。\src{REG-009}\src{REG-010}\src{REG-012}\src{REG-007}\src{REG-008}

\clearpage
\section{后续验证：以模型任务和产物确定原型}
同行案例已提供可以借鉴的模型任务，但本地部署的价值仍需要在实际数据和环境中验证。初期宜选择输入范围明确、输出可以独立核对的任务，再逐步增加工具和自主执行能力。任务要求由此成为模型选型和资源配置的依据。

\begin{table}[H]\small\centering
\caption{可用于本地验证的LLM任务}
\begin{tabularx}{\linewidth}{@{}P{26mm}Y Y Y@{}}\toprule
任务 & 提供给模型的输入 & 预期模型产物 & 独立核验方式\\\midrule
事件提取 & 授权公告、字段定义和输出格式。 & 事件、状态、数量、单位及原文证据。 & 人工标签核对，程序验证格式与时间。\\
只读查询 & 问题、数据字典或接口目录。 & 查询参数或SQL，以及结果答复。 & 人工标准答案、执行权限和口径检查。\\
资料问答 & 问题、获准检索片段和来源。 & 带引用的答案或资料不足说明。 & 检查原文支持、遗漏及权限隔离。\\
因子规格或代码 & 研究要求、可用变量、算子及接口定义。 & 因子表达式、实现或修正建议。 & 程序执行、语义核验和时序检查，再做研究评价。\\\bottomrule
\end{tabularx}
\end{table}

按项目计划，下一阶段比较5--8个开放模型，并在本地GPU上部署一个模型完成金融文本任务。开发阶段使用开发样本调整提示词和工具，测试阶段保持数据与规则固定。除最终准确率，还应记录结构化输出解析失败、工具选择错误、无依据生成和修复失败等问题，避免只报告成功示例。

效率评价应覆盖模型请求、工具执行和人工复核的总耗时。成本评价则在后续三年TCO分析中纳入设备、API、工程接入、运行维护和错误修正。当前报告仅形成模型任务与验证条件，不将公开案例效果视为本项目已实现的效率、收益或合规结论。
\clearpage

\appendix
\section{公司规模、模型任务与技术合作信息}
表A汇总20家主要样本与大成补充案例。规模为2026年二季度末非货公募规模，剔除ETF联接基金，单位为亿元人民币。任务与产物仅按公开内容填写，尚未明确LLM分工的材料保留相应说明。\src{FUND-400}
{\footnotesize
\begin{longtable}{@{}P{17mm}P{20mm}P{35mm}P{38mm}P{41mm}@{}}
\caption{公司结构化数据库摘要：规模单位为亿元人民币；规模数值统一来自FUND-400。}\label{tab:landscape}\\
\toprule 公司 & 2026Q2规模 & 场景及LLM边界 & 披露技术 / 工具 & 合作方与证据状态\\\midrule\endfirsthead
\multicolumn{5}{l}{\small 表\ref{tab:landscape}（续）}\\\toprule 公司 & 2026Q2规模 & 场景及LLM边界 & 披露技术 / 工具 & 合作方与证据状态\\\midrule\endhead
\midrule\multicolumn{5}{r}{\scriptsize 接下页}\endfoot
\bottomrule\endlastfoot
易方达 & 17,029.86 & 指数研究、基金研选；论文、比赛与生产系统分列 & 指数一体化、Index-Hub客户端；模型后端未完整公开 & 合作研究包含高校；公开材料未完整揭示生产供应商职责。\src{FUND-004}\src{FUND-115}\\
华夏 & 11,987.79 & 模型安全接入；采购不能证明验收上线 & 大模型安全网关；2025采购证据 & 北京世纪东凌中标；交付范围和验收待核。\src{FUND-001}\\
广发 & 10,355.25 & 智能理财；2026调研称LLM挖掘另类因子 & 协议列DeepSeek/Qwen等模型家族；量化实现未详述 & 模型来源不等于实施商；合作合同未公开。\src{FUND-009}\src{FUND-251}\\
富国 & 9,889.56 & 市场观点、本地知识问答、政策文本因子 & LLM/RAG；Ray、ClickHouse等平台技术 & 自研案例；外部承包及分工未披露。\src{FUND-007}\\
南方 & 7,856.35 & 交易助理：受控SQL、API、RAG、Python & 意图解析、任务分发；选表与执行校验 & 案例未明确运行模型及外部技术承包方。\src{FUND-200}\\
汇添富 & 7,719.67 & 托管风险邮件提取、条款问答；规则引擎另列 & DeepSeek-R1、QianWen、FastAPI；Flink/Doris & 模型与组件名称已见图示；服务商及合同职责未披露。\src{FUND-220}\\
景顺长城 & 7,465.45 & 报道本地LLM文档解析与检索 & DeepSeek、OCR、Embedding、重排（报道） & 公司原始技术实施材料存在缺口；供应商待核。\src{FUND-251}\\
嘉实 & 6,911.96 & ESG数据与研究；已归档AI/NLP不全属LLM & 可持续投资报告中的AI/NLP & 有公司原件；具体LLM流程及合作分工未确认。\src{FUND-304}\\
博时 & 6,386.18 & BSBox研究工具、数据访问与任务协作 & 内部模型、Skills、用户容器、审计及记忆 & 2026报道；未取得完整量化研究/回测任务轨迹。\src{FUND-215}\\
鹏华 & 6,366.21 & 投研资讯提取、知识问答、营销合规 & 私有多模态模型、Prompt治理、RAG、PSG网关 & 公司原始案例；外部实施方未披露。\src{FUND-308}\\
招商 & 6,276.62 & 个股洞察、财报点评；量化回测为探索方向 & 多工作流与提取模型（媒体调研） & 公司原始技术稿有缺口；不把传统量化工具算作LLM。\src{FUND-251}\\
国泰 & 6,238.11 & 组合管理、检索与问答；传统NER单列 & Ascend Atlas 910B、DeepSeek/RAG；BERT等另列 & 公司原始平台案例；外部承包分工未披露。\src{FUND-403}\\
永赢 & 5,332.53 & 安全GPT识别钓鱼邮件；属运营风险 & 大模型邮件安全处理 & 2026协会原始案例；非量化研究系统。\src{FUND-352}\\
工银瑞信 & 5,006.18 & 养老金运营平台、FundGPT；与传统因子模型区分 & 大模型平台及相关合作披露 & 历史供应商材料涉及恒生聚源、工银科技、智谱AI；现行合同和具体分工待核。\src{FUND-312}\src{FUND-316}\\
天弘 & 4,824.03 & FinAgent财务推理、检索及工具协作 & Think、THAI、弘思、智搜；Wind/同花顺/MCP & 数据工具不等于实施供应商；承包职责未公开。\src{FUND-300}\\
中欧 & 4,687.70 & 数据抓取、财务建模、报告写作Skills & Agent与100余Skills（披露） & 技术原件仍有缺口；不能把报道中的其他公司案例归属中欧。\src{FUND-356}\src{FUND-251}\\
华安 & 4,632.07 & 灵思研究/知识平台；2025历史补充 & 本地Qwen/DeepSeek-R1与外部API、SFT & 有原始案例集；不称2026新项目，实施方未详述。\src{FUND-150}\\
华泰柏瑞 & 4,459.62 & 文本结构化、事件与风险跟踪；自主回测Agent属展望 & 文本处理与量化信号融合；模型/Prompt未公开 & 已恢复2026微信原刊，仍为记者叙述，缺公司技术原件。\src{FUND-359}\\
兴证全球 & 3,751.87 & 千询/兴宝固收交易；NLP与LLM边界需逐模块核实 & 交易信息抽取及QTrade/iDeal等工具衔接 & 有公司平台案例；基础LLM型号与技术承包角色未明确。\src{FUND-011}\\
平安 & 3,447.69 & AI青蚨：文档、检索与投研材料组织 & 资料检索、可追溯输出；具体模型未披露 & 2026媒体材料；公司原始技术实施证据仍不足。\src{FUND-160}\\
大成\newline（补充） & 不纳入同口径排名 & 固收全链路中的交易要素提取 & Qwen2.5；Flink、SpringCloud、Vue、O32衔接 & 公司原始案例；外部合作方未具名。\src{FUND-013}\\
\end{longtable}}
\smallnote{五家核心样本的一手实施证据缺口：景顺长城、招商、中欧、华泰柏瑞、平安。其余公司的“一手证据”可能仅覆盖历史AI模块、协议或采购，不等于2026 LLM生产应用已独立核实。中信建投基金等新增机构见知识库扩展资料，不并入这21家公司计数。}

\clearpage
\section{相关法律、监管文件与行业标准}
以下文件提供条款查阅位置，资料截止2026年9月9日。实际适用情况仍需结合机构、数据和用途判断。
{\fontsize{9.5}{12}\selectfont\renewcommand{\arraystretch}{1.13}
\begin{longtable}{@{}P{43mm}P{37mm}P{72mm}@{}}
\toprule 文件及性质 & 定位 & 对本项目的适用要点\\\midrule\endfirsthead
\toprule 文件及性质 & 定位 & 对本项目的适用要点\\\midrule\endhead
\bottomrule\endlastfoot
生成式人工智能服务管理暂行办法（部门规章）\src{REG-001} & 第2、4(5)、17条 & 先判断是否向境内公众提供服务；透明、准确可靠及备案／评估义务须结合适用范围和条件。内部应用排除条款不免除其他法律责任。\\
个人信息保护法（法律）\src{REG-002} & 第13、21、24、27、54--56条 & 多种合法基础；委托约定与监督；自动化决策的透明、公平和特定说明／拒绝权；公开信息合理范围；定期审计及特定处理事前影响评估。评估报告和处理记录至少保存3年。\\
证券基金经营机构信息技术管理办法（证监会规章，2021修订）\src{REG-003} & 第3、43--45、63条 & 按实际IT服务角色判断；保留自身责任和重要系统控制，外部独立运维限制有法定例外；审查报送、替换预案及协议。第63条定义重要信息系统。\\
网络数据安全管理条例（行政法规，2025施行）\src{REG-004} & 第9、12、19、29、31、33、37条 & 安全措施、委托／提供个人信息与重要数据的合同监督及处理记录；训练数据规则结合生成服务角色理解；重要数据按目录、告知及识别义务判断。\\
证券期货业网络和信息安全管理办法（令218）\src{REG-005} & 第2、4、29--34、71、75条 & 证券基金机构网络信息安全及投资者信息保护责任；第75条废止旧信息安全保障办法。\\
促进和规范数据跨境流动规定（部门规章）\src{REG-006} & 第5(4)、7--8、10--11、13条 & 出境机制按关基／重要数据和累计人数等分层；普通与敏感个人信息门槛不同；豁免机制仍保留合法处理等义务。需公司级事实，不能按原型样本量代替累计统计。\\
网络安全法（2025修正，2026施行）\src{REG-013} & 第23条 & 等保相关措施、网络运行和安全事件监测记录；不少于6个月的期限针对相关网络日志，不能直接等同全部提示词原文保留期。\\
数据安全法（法律）\src{REG-014} & 第21、27、29--32条 & 数据分类分级、全流程管理、风险监测、重要数据风险评估和合法收集。保护对象不限于个人信息。\\
个人信息保护合规审计管理办法（2025施行）\src{REG-015} & 第3--4条 & 定期开展合规审计；处理超过1,000万人个人信息的，至少每两年一次。低于门槛不等于免审计。\\
人工智能生成合成内容标识办法（2025施行）\src{REG-016} & 第2--6、9条 & 按生成服务、传播平台、发布者及具体内容类型判断显式／隐式标识与用户声明义务；对外发布需重判角色。配套强制国标技术实现需在架构确定后核对。\\
基金经营机构大模型技术应用规范（团体标准）\src{REG-009}\src{REG-010} & 9.1.2、9.2、11.2--11.3；附录A／B & 自有场景评测、部署、数据权限与解释参考；11.3(e)用“宜”，附录为资料性。适用依据还须结合标准化法第18条与实际采用约定。\src{REG-012}\\
SR 26-2及配套银行MRM指引（美国监督指引）\src{REG-007}\src{REG-008} & 发布函；指引I--II节、脚注3 & 2026-04-17替代旧指引；非可强制执行标准，面向所述银行机构。生成式／Agentic AI不在定义范围，机构仍自行确定风险治理。非MSIM内部政策。\\
\end{longtable}}

\clearpage
\section{研究结果与实验条件补充}
以下保留主要论文和机构案例的结果，用于理解方法验证的范围。结果均由原作者报告，未在本项目中独立复现，不代表基金实盘表现。
\setcounter{table}{0}
\renewcommand{\thetable}{C.\arabic{table}}
\subsection{MENTOR行业排序}
模型输入文本覆盖2023--2024年，包括160,190篇中文KOL帖子与72,108篇英文财经新闻，A股和美股行情分别来自Wind与Bloomberg。排序评价涉及9个与11个行业。\src{FUND-113}\src{FUND-114}

\begin{table}[H]\small\centering
\caption{MENTOR论文中的行业排序结果}\label{tab:mentor}
\begin{tabularx}{\linewidth}{@{}P{41mm}Y Y Y@{}}\toprule
方法 & 美股行业排序相关性 & A股行业排序相关性 & A股前3行业命中率\\\midrule
动量基线 & 0.159 & 0.100 & 0.444\\
SEP + DeepSeek & 0.078 & 0.118 & 0.472\\
MENTOR + DeepSeek & 0.220 & 0.141 & 0.488\\
MENTOR + O1 & 0.221 & 0.173 & 0.439\\\bottomrule
\end{tabularx}
\source{MENTOR正文表3（PDF第11页）。排序相关性采用Spearman系数；前3命中率衡量预测与实际领先行业的重合程度。数字为原作者报告值。}
\end{table}
表中MENTOR改善了部分排序相关性，但A股O1版本的前3命中率低于DeepSeek版本与动量。完整版权数据未公开，正文和补充材料部分数据描述对应不清，预训练污染与样本时间仍需复核。

\subsection{Man AlphaTrend验证实验}
作者故意从既有突破信号中移除提高反应速度的机制，再输入预期有效、预期较差和开放性想法，检查研究流程能否形成相应结论。\src{FUND-323}

\begin{table}[H]\small\centering
\caption{AlphaTrend对三类研究想法的检验}\label{tab:man}
\begin{tabularx}{\linewidth}{@{}P{38mm}Y Y@{}}\toprule
研究命题 & 原作者观察到的结果 & 结果的含义\\\midrule
恢复被移除的反应机制 & 方案的Sharpe集中在约1.05，基准约为0.95。 & 系统识别出预先设定的改进机会。\\
改用局部峰谷等信号 & Sharpe中心约为0.80，几乎所有方案落后于基准。 & 系统能够给出负面结果，帮助停止无效尝试。\\
在更宽的范围探索趋势方案 & Sharpe约为0.70--1.00，中心约0.85，不同时期表现不同。 & 结果存在条件差异，需要进一步限定假设与适用范围。\\\bottomrule
\end{tabularx}
\source{Man AHL，2026-02-11，图2--4及相邻正文。结果是原作者的历史模拟展示，不代表真实产品业绩。}
\end{table}
恢复已知机制主要验证流程能否工作，不代表发现新Alpha。原文未完整披露资产池、参数、成本和统计检验；图注截至2015年，正文部分表述涉及2012--2016年，不能补写统一样本区间。同一开放想法下，作者报告Claude 4.0 Sonnet方案的大部分两两相关性高于0.85，GPT-5约为0.75--1.0，仅说明该设置中的候选差异。

\subsection{QuantaAlpha总体实验及案例口径}
表中使用同一GPT-5.2模型比较三个LLM框架，各方法产生约150个通过验证的因子，再交给同一LightGBM模型评价。\src{EXT-QA}\src{EXT-QC}

\begin{table}[H]\small\centering
\caption{QuantaAlpha论文中的同模型对照}\label{tab:qa}
\begin{tabularx}{\linewidth}{@{}P{43mm}Y Y Y@{}}\toprule
框架（均用GPT-5.2） & IC & 年化超额收益 & 超额收益最大回撤\\\midrule
RD-Agent & 0.0286 & 3.58\% & 16.76\%\\
AlphaAgent & 0.0347 & 1.11\% & 13.89\%\\
QuantaAlpha & 0.0472 & 4.68\% & 11.80\%\\\bottomrule
\end{tabularx}
\source{QuantaAlpha v3表1及附录A--B。IC衡量预测值与后续收益的相关性，表中数字均为原作者报告值。}
\end{table}
总体市场为沪深300，训练期2016--2020年，验证期2021年，测试期2022年初至2025年12月26日。组合每日选择50只等权股票，每次替换5只，采用次日开盘价，买入与卖出费率分别为0.05\%和0.15\%。收益与回撤基于扣费后的超额收益；预测标签为下一交易日到再下一交易日的收盘收益。其他传统或深度方法在部分组合指标上更好，不能据该节选认定LLM全面占优。

案例表达式的价格与成交量变动均以当日值作分母，日内收益也以当日收盘价作分母。复现时应按原表达式保留这一口径。

使用DeepSeek-V3.2的消融实验中，移除修改环节后，IC由0.0461降至0.0382，年化超额收益由4.53\%降至3.27\%。附录C案例则报告子因子Rank IC为0.0311，父记录为0.0216和0.0246，信息比率由0.973降至0.963，超额收益最大回撤由约7.30\%扩大至11.37\%。该案例与总体实验的轮次、换手及成本口径不完全一致，正文仅用其说明模型解释反馈和提出修改的过程，不合并为总体收益证据。

作者报告一次主要运行约20小时、约180万tokens，采用托管模型API与CPU回测。公开代码为选取文件存档，配置、数据许可和完整运行环境仍需核对，尚未执行复现。

\subsection{AlphaQT-Bench代码生成评价}
该基准包含270项任务并测试12个模型，LLM返回Python实现，测试程序分别检查运行、时序因果、功能和计算结构。四项同时通过的平均验证正确率（Verified Accuracy），在论文模型样本中最高为93.7\%（Gemini-2.5-Pro），开放权重模型最高为81.9\%（DeepSeek-V3）。该结果评价任务实现，不等于策略成功率或本地推理性能。\src{EXT-AB}

\clearpage
\section{参考文献}
文献按正文与附录首次引用的顺序排列。未特别说明时，页码指原始PDF物理页。来源原件及固定代码版本的记录另附。
\referenceentry{FUND-007}{『基金行业金融科技发展奖』富国基金：指数基金智能投资决策系统}{2026-01-19协会目录；PDF未注明日期}{PDF第8–9页；本地RAG、市场观点与因子辅助评价}{https://www.amac.org.cn/xwfb/hydt/202601/P020260202337808067241.pdf}

\referenceentry{FUND-200}{『基金行业金融科技获奖成果宣传活动』南方基金：基于多智能体协同的交易助理}{2026-01-23协会目录；PDF未注明修订日期}{PDF第4–8页，尤其第5–6页NL2SQL及权限检查}{https://www.amac.org.cn/xwfb/hydt/202601/P020260202337849876845.pdf}

\referenceentry{FUND-308}{『基金行业金融科技获奖成果宣传活动』鹏华基金：基于动态思维链的资管业务智能体建设}{2026-01-27协会目录；PDF未注明日期}{PDF第3–6页；Prompt治理、RAG与营销材料审核}{https://www.amac.org.cn/xwfb/hydt/202601/P020260202337850951100.pdf}

\referenceentry{FUND-323}{A Trend Following Deep Dive: AlphaTrend and Agentic Research Workflows}{2026-02-11}{Man AHL；原文流程图与三组研究想法实验}{https://www.man.com/insights/alphatrend-agentic-research-workflows}

\referenceentry{EXT-QA}{QuantaAlpha: An Evolutionary Framework for LLM-Driven Alpha Mining}{2026；论文封面日期2026-05-19（arXiv v3）}{v3；表1--2、附录A--C及方法；作者团队含高校与研究项目}{https://arxiv.org/abs/2602.07085v3}

\referenceentry{FUND-400}{公募非货规模最新座次出炉，万亿公募增至3家，中段座次重排}{2026-07-21}{财联社；2026Q2非货公募规模，剔除ETF联接基金}{https://www.cls.cn/detail/2432695}

\referenceentry{FUND-300}{『基金行业金融科技发展奖』天弘基金：基于大模型的 FinAgent 金融智能体系统}{2026-01-19协会目录；PDF未注明日期}{PDF第4–7页；FinAgent分层与业务拆解}{https://www.amac.org.cn/xwfb/hydt/202601/P020260202337562370809.pdf}

\referenceentry{FUND-220}{『基金行业金融科技获奖成果宣传活动』汇添富基金：智汇投资风险管理平台}{2026-01-27协会目录；PDF未注明修订日期}{PDF第5–11页，尤其第7页图2；LLM与规则平台边界}{https://www.amac.org.cn/xwfb/hydt/202601/P020260202337837900594.pdf}

\referenceentry{FUND-215}{智启金融新范式，博时前瞻布局企业级“龙虾”生态，赋能资管全链条}{2026-04-30}{博时BSBox报道；工具与运行环境}{https://finance.sina.com.cn/money/fund/jjh/2026-04-30/doc-inhwfyep1330822.shtml}

\referenceentry{FUND-356}{中欧基金窦玉明：AI赋能时代，「三化」协同夯实长期业绩根基}{2026-06-01}{中欧公众号转载中国基金报；Skills与投研协作}{https://mp.weixin.qq.com/s/rgMRw8HbZUhiu9PkUBM7vA}

\referenceentry{FUND-251}{公募基金AI部署走到哪一步了？20家头部公募AI部署全景调研}{2026-08-20}{中国基金报20家调研；媒体/公司归属陈述，不替代技术原件}{https://www.chnfund.com/article/AR99eb77fa-e926-d40e-abb2-3a232d36f0d0}

\referenceentry{FUND-359}{华泰柏瑞量化团队：“人机协同”迈入深水区，AI投研构筑长期护城河}{2026-09-03}{中国证券报原刊；文本处理实践与Agent展望分开}{https://mp.weixin.qq.com/s?__biz=MjM5MzMwNjM0MA\%3D\%3D&idx=2&mid=2651361688&sn=465ee3d543763b18d8ee971802cc5fcb}

\referenceentry{FUND-001}{关于华夏基金管理有限公司采购结果的公告}{2025-12-31；历史采购证据}{采购公告；不证明验收或完整LLM部署}{https://www.chinaamc.com/upload/resources/file/2025/12/31/442744.pdf}

\referenceentry{FUND-316}{恒生电子2023可持续发展报告}{2023报告期；历史合作证据}{恒生电子可持续发展报告；FundGPT相关合作方}{https://static.cninfo.com.cn/finalpage/2024-03-25/1219392176.PDF}

\referenceentry{FUND-113}{MENTOR: a multi-agent framework for event and narrative trend prediction with optimized reasoning}{2025，26(10):1847–1861}{PDF第8–11页及表3；方法与行业排序结果}{https://doi.org/10.1631/FITEE.2500608}

\referenceentry{FUND-114}{MENTOR supplementary materials}{2025；MENTOR官方补充材料}{PDF第1–5页，尤其第2–3页Prompt}{https://jzus.zju.edu.cn/downloadesm.php?filepath=.\%2Fdownload\%2FESM\%2FFITEE\%2FESM2500608.pdf&filetype=pdf}

\referenceentry{FUND-115}{指数直通车 Index-Hub Skills：官方查询客户端与接口文档}{代码快照；不推断发布日期}{固定提交1058e1b；公开查询客户端与配置缺口}{https://github.com/EFundAI/index-hub-skills/tree/1058e1bfdaaf3dc1ee025ba2d2e19ea000657fd4}

\referenceentry{FUND-116}{指数直通车 AI Skills 帮助文档}{无日期帮助文档；2026-09-09归档}{PDF第8–10页查询示例；非服务端代码}{https://cdn.efunds.com.cn/eda/h5/itcenter/pd/ai-skills-doc/help.pdf}

\referenceentry{REG-010}{基金经营机构大模型技术应用规范 — T/AMAC 0004-2026}{2026-04-03}{适用条件与精确条款见附录B}{https://www.amac.org.cn/xwfb/xhyw/202604/P020260429638714694502.pdf}

\referenceentry{FUND-119}{NLPCC 2026 Shared Task 4: Investment Agent Starter}{NLPCC 2026共享任务；仓库快照}{固定提交4744a77；研究基线、Prompt、数据及回测客户端}{https://github.com/splash-li/NLPCC2026-Shared-Task-4/tree/4744a772fda59ef886c525cf3d4b4ac3c992710f}

\referenceentry{EXT-QC}{QuantaAlpha官方研究实现（固定版本）}{固定提交 b7ceb27b1001}{方法实现；代码固定版本见原始链接；未执行复现}{https://github.com/QuantaAlpha/QuantaAlpha/tree/b7ceb27b1001261d7a95b209a963664ae1f8ab23}

\referenceentry{EXT-AB}{AlphaQT-Bench：量化因子代码评价基准}{ACL 2026 Findings}{16页论文；第16页Prompt及执行、因果、语义、结构评价}{https://aclanthology.org/2026.findings-acl.138/}

\referenceentry{REG-001}{生成式人工智能服务管理暂行办法}{2023-07-13； 签署 2023-07-10}{适用条件与精确条款见附录B}{https://www.cac.gov.cn/2023-07/13/c_1690898327029107.htm}

\referenceentry{REG-005}{证券期货业网络和信息安全管理办法}{令文日期 2023-02-27； 发布 2023-03-03}{适用条件与精确条款见附录B}{https://www.csrc.gov.cn/csrc/c101953/c7202800/7202800/files/\%E9\%99\%84\%E4\%BB\%B61\%EF\%BC\%9A\%E8\%AF\%81\%E5\%88\%B8\%E6\%9C\%9F\%E8\%B4\%A7\%E4\%B8\%9A\%E7\%BD\%91\%E7\%BB\%9C\%E5\%92\%8C\%E4\%BF\%A1\%E6\%81\%AF\%E5\%AE\%89\%E5\%85\%A8\%E7\%AE\%A1\%E7\%90\%86\%E5\%8A\%9E\%E6\%B3\%95.pdf}

\referenceentry{REG-002}{中华人民共和国个人信息保护法}{2021-08-20}{适用条件与精确条款见附录B}{https://www.cac.gov.cn/2021-08/20/c_1631050028355286.htm}

\referenceentry{REG-003}{证券基金经营机构信息技术管理办法}{2018-12-19； 修订 2021-01-15}{适用条件与精确条款见附录B}{https://www.csrc.gov.cn/csrc/c106256/c1653896/content.shtml}

\referenceentry{REG-004}{网络数据安全管理条例}{令文日期 2024-09-24； 网信办发布 2024-09-30}{适用条件与精确条款见附录B}{https://www.cac.gov.cn/2024-09/30/c_1729384452307680.htm}

\referenceentry{REG-014}{中华人民共和国数据安全法}{2021-06-10}{适用条件与精确条款见附录B}{https://www.cac.gov.cn/2021-06/11/c_1624994566919140.htm}

\referenceentry{REG-006}{促进和规范数据跨境流动规定}{2024-03-22}{适用条件与精确条款见附录B}{https://www.cac.gov.cn/2024-03/22/c_1712776611775634.htm}

\referenceentry{REG-009}{关于发布《基金经营机构大模型技术应用规范》团体标准的公告}{2026-04-03}{适用条件与精确条款见附录B}{https://www.amac.org.cn/xwfb/xhyw/202604/t20260403_27458.html}

\referenceentry{REG-012}{中华人民共和国标准化法（2017年修订）}{修订 2017-11-04； 国家知识产权局转载 2019-07-29}{适用条件与精确条款见附录B}{https://www.cnipa.gov.cn/art/2019/7/29/art_104_67809.html}

\referenceentry{REG-007}{SR 26-2: Revised Guidance on Model Risk Management}{2026-04-17}{适用条件与精确条款见附录B}{https://www.federalreserve.gov/supervisionreg/srletters/SR2602.htm}

\referenceentry{REG-008}{Supervisory Guidance on Model Risk Management}{2026-04-17}{适用条件与精确条款见附录B}{https://www.federalreserve.gov/frrs/guidance/supervisory-guidance-on-model-risk-management.htm}

\referenceentry{FUND-004}{『基金行业金融科技获奖成果宣传活动』易方达基金：基于数智赋能的指数业务一体化平台}{2026年1月协会案例；PDF未注明修订日期}{指数业务一体化平台；平台指标与LLM效果须分开}{https://www.amac.org.cn/xwfb/hydt/202601/P020260202337819827722.pdf}

\referenceentry{FUND-009}{广发基金智能理财助理服务协议}{协议未注明发布日期}{大模型技术名称与服务协议；不据URL月份认定发布日期}{https://cdnwww.gffunds.com.cn/gfjjnew/jjgg/lsgg/202603/W020260327646002876113.pdf}

\referenceentry{FUND-304}{嘉实基金2025可持续投资报告}{2025报告期；发布日期未独立确认}{嘉实可持续投资报告；AI/NLP不自动等于LLM}{https://www.jsfund.cn/ueditor/jsp/upload/file/20260331/1774939300245008932.pdf}

\referenceentry{FUND-403}{『基金行业金融科技获奖成果宣传活动』国泰基金：基金组合管理平台建设}{2026年1月协会案例；PDF未注明日期}{国泰组合平台；DeepSeek与传统NER模块分开}{https://www.amac.org.cn/xwfb/hydt/202601/P020260202337808791393.pdf}

\referenceentry{FUND-352}{AI筑牢邮件安全防线 永赢基金以大模型守护金融数字资产}{2026-07-01协会目录；PDF未注明日期}{永赢安全GPT邮件反钓鱼；非投研量化系统}{https://www.amac.org.cn/xwfb/hydt/202607/P020260701319545695145.pdf}

\referenceentry{FUND-312}{『基金行业金融科技获奖成果宣传活动』工银瑞信基金：基于人工智能的一站式养老金投资运营管理平台}{2026-01-21协会目录；PDF未注明日期}{工银瑞信养老金运营管理平台}{https://www.amac.org.cn/xwfb/hydt/202601/P020260202337848496316.pdf}

\referenceentry{FUND-150}{2025 CCF企业数字化发展推荐案例集 — 华安基金管理有限公司——华安“灵思”AI智能创新平台}{2025版案例集；历史补充}{华安“灵思”；本地与API混合平台}{https://edc-book.ccf.org.cn/\%E3\%80\%8A2025CCF\%E4\%BC\%81\%E4\%B8\%9A\%E6\%95\%B0\%E5\%AD\%97\%E5\%8C\%96\%E5\%8F\%91\%E5\%B1\%95\%E6\%8E\%A8\%E8\%8D\%90\%E6\%A1\%88\%E4\%BE\%8B\%E9\%9B\%86\%E3\%80\%8B.pdf}

\referenceentry{FUND-011}{『基金行业金融科技获奖成果宣传活动』兴证全球基金：千询固收智能交易平台}{2026年1月协会案例；PDF未注明日期}{千询固收交易平台；NLP与基础LLM需区分}{https://www.amac.org.cn/xwfb/hydt/202601/P020260202337841539992.pdf}

\referenceentry{FUND-160}{平安基金发布“AI青蚨”：破解数据“沉睡”困局 构建智能投研新范式}{2026-02-06}{中国证券报；平安AI青蚨报道}{https://www.cs.com.cn/ssgs/gsxl/202602/t20260206_6537131.html}

\referenceentry{FUND-013}{『基金行业金融科技获奖成果宣传活动』大成基金：固收全链路数智一体化平台}{2026年1月协会案例；PDF未注明日期}{固收平台中的Qwen2.5交易要素提取}{https://www.amac.org.cn/xwfb/hydt/202601/P020260202337853082516.pdf}

\referenceentry{REG-013}{中华人民共和国网络安全法（2025年修正）}{2025-10-28}{适用条件与精确条款见附录B}{https://www.cac.gov.cn/2025-12/29/c_1768735112911946.htm}

\referenceentry{REG-015}{个人信息保护合规审计管理办法}{2025-02-12}{适用条件与精确条款见附录B}{https://www.cac.gov.cn/2025-02/14/c_1741233507681519.htm}

\referenceentry{REG-016}{人工智能生成合成内容标识办法}{2025-03-07}{适用条件与精确条款见附录B}{https://www.cac.gov.cn/2025-03/14/c_1743654684782215.htm}

\end{document}
